% ============ §31 Chains of prime ideals (原书 pp.246-254) ============
\section{Chains of prime ideals}\label{sec:31}\index{chains of prime ideals, section 31}\index{catenary ring, section 31}\index{equidimensional ring, section 31}\index{formally equidimensional, section 31}\index{formally catenary, section 31}\index{universally catenary, section 31}

\begin{theorem}\label{thm:31.1}
Let $A$ be a Noetherian ring and $P \in \Spec(A)$. Then there are at most finitely
many prime ideals $P'$ of $A$ satisfying $P \subset P'$, $\htop(P'/P) = 1$ and
$\htop P' > \htop P + 1$. (Ratliff [3] in the semi-local case, and McAdam [3] in
the general case.)
\end{theorem}

\begin{proof}
Let $\htop P = n$ and take $a_1, \ldots, a_n \in P$ such that $\htop(a_1, \ldots, a_n) = n$. Set
$I = (a_1, \ldots, a_n)$ and let $P_1 = P, P_2, \ldots, P_r$ be the minimal prime divisors of $I$.
In general, if $\{Q_\lambda\}_\Lambda$ is an infinite set of prime ideals such that
$Q_\lambda \supset P$ and $\htop(Q_\lambda/P) = 1$, then $\bigcap_\Lambda Q_\lambda = P$. This is because
$\bigcap Q_\lambda$ is equal to its own radical, hence is a finite intersection of prime
ideals containing $P$, hence either
\[
  \bigcap_\Lambda Q_\lambda = P
\]
or
\[
  \bigcap_\Lambda Q_\lambda = Q_{\lambda_1} \cap \cdots \cap Q_{\lambda_k};
\]
but the second case cannot occur since $Q_{\lambda_j} \not\supset Q_{\lambda_j}$ for
$j \ne \{\lambda_1, \ldots, \lambda_k\}$. Therefore if there were an infinite number of prime
ideals $Q_\lambda$ such that $Q_\lambda \supset P$, $\htop(Q_\lambda/P) = 1$ and
$\htop Q_\lambda \ne n + 1$, then $\bigcap Q_\lambda = P$. Hence there would exist a
$Q_\lambda$, say $Q_0$, which does not contain $P_1, \ldots, P_r$. Then $P$ would be the only
prime ideal lying between $Q_0$ and $I$. Let $b \in Q_0 - P$. Then $Q_0$ would be a
minimal prime divisor of $I + bA = (a_1, \ldots, a_n, b)$, so that
$\htop Q_0 \le n + 1$, a contradiction.
\end{proof}

\begin{theorem}[\textbf{Ratliff's weak existence theorem}]\label{thm:31.2}
Let $A$ be a Noetherian ring, and $p$, $P \in \Spec A$ be such that $p \subset P$,
$\htop p = h$ and $\htop(P/p) = d > 1$; then there exist infinitely many $p' \in \Spec A$ with
the properties
\[
  p \subset p' \subset P, \quad \htop p' = h + 1 \quad \text{and} \quad \htop(P/p') = d - 1.
\]
\end{theorem}

\begin{proof}
We first observe that if $P \supset p_1 \supset p_2 \supset \cdots \supset p_d = p$ is a strictly
decreasing chain of prime ideals, and if $p_{d-2} \subset p' \subset p$ with
$\htop p' = h + 1$ then $\htop(P/p') = d - 1$. Now there exist infinitely many
$p' \in \Spec A$ such that $p_{d-2} \subset p' \subset p$ and $\htop(p'/p) = 1$: for if
$p_1', \ldots, p_k'$ are a finite set of these, let $a \in p_{d-2} - \bigcup p_i'$, and let
$p_{k+1}'$ be a minimal prime divisor of $p + aA$ contained in $p_{d-2}$; then
$\htop(p_{k+1}'/p) = 1$. By the previous theorem, all but finitely many of these
satisfy $\htop p' = \htop p + 1$.
\end{proof}

\begin{lemma}\label{lem:31.1}
Let $A$ be a Noetherian ring, and $P \in \Spec A$ with $\htop P = h > 1$; suppose that
$u \in P$ is such that $\htop(uA) = 1$. Then there exist infinitely many prime ideals
$Q \subset P$ such that
\[
  u \notin Q \quad \text{and} \quad \htop Q = h - 1.
\]
\end{lemma}

\begin{proof}
Suppose that $p_1, \ldots, p_t$ are the minimal prime ideals of $A$, and let
$P_1, \ldots, P_s$ be finitely many given height $1$ prime ideals not containing $u$.
Let $Q_1, \ldots, Q_m$ be the minimal prime divisors of $uA$, so that these are also
height $1$ prime ideals of $A$. Since $h > 1$, there exists $v \in P$ not contained in
any $p_i$, $P_j$ or $Q_k$. Then
\[
  \htop(u, v) = 2 \quad \text{and} \quad \htop(v) = 1.
\]
Now let $P_{s+1}, \ldots, P_{s'}$ be the minimal prime divisors of $(v)$; continuing in the
same way we can find infinitely many height $1$ prime ideals not containing $u$, so
that if $h = 2$ we are done. If $h > 2$ we set $\bar{A} = A/(u)$ and
$\bar{P} = P/(u)$, so that $\htop \bar{P} = h - 1$, and since the image $\bar{v}$ of $v$ in
$\bar{A}$ satisfies $\htop(\bar{v}) = 1$, by induction on $h$ we can find infinitely many prime
ideals $\bar{Q}$ of $\bar{A}$ satisfying
\[
  \bar{Q} \subset \bar{P} \quad \text{and} \quad \htop \bar{Q} = h - 2.
\]
The inverse image $Q$ of $\bar{Q}$ in $A$ does not contain $u$, and from
$Q/(u) = \bar{Q}$ we have $\htop Q = h - 1$.
\end{proof}

\begin{theorem}[\textbf{Ratliff's strong existence theorem}]\label{thm:31.3}
Let $A$ be a Noetherian integral domain, $p$, $P \in \Spec A$, and suppose that
$\htop p = h > 0$ and $\htop(P/p) = d$. Then for each $i$ with $0 \le i \le d$ the set
\[
  \{p' \in \Spec A \mid p' \subset P, \htop(P/p') = d - i \text{ and } \htop p' = h + i\}
\]
is infinite.
\end{theorem}

\begin{proof}
For $i > 0$ this follows at once from the weak existence theorem, so that we
consider the case $i = 0$.

\textit{Step 1.} Replacing $A$ by $A_p$ we can assume that $(A, P)$ is a local integral
domain. Choose $a_1, \ldots, a_h$ such that $\htop(a_1, \ldots, a_j) = j$ for $j = 1, \ldots, h$,
and set
\[
  \mathfrak{a} = (a_1, \ldots, a_h) \quad \text{and} \quad \mathfrak{b} = (a_1, \ldots, a_{h-1}).
\]
Then $p$ is a minimal prime divisor of $\mathfrak{a}$. Let
\[
  \mathfrak{a} = \mathfrak{a}_1 \cap \cdots \cap \mathfrak{a}_r \quad \text{and} \quad
  \mathfrak{b} = \mathfrak{b}_1 \cap \cdots \cap \mathfrak{b}_s
\]
be shortest primary decompositions of $\mathfrak{a}$ and $\mathfrak{b}$, and let $p_i$, $p_j'$ be the
prime divisors of $\mathfrak{a}_i$ and $\mathfrak{b}_j$, respectively; we can assume that
$p = p_1$. Suppose that $\mathfrak{b}_{t+1}, \ldots, \mathfrak{b}_s$ are all the $\mathfrak{b}_j$ not contained
in $p$. Then since
\[
  \mathfrak{a}_2 \cap \cdots \cap \mathfrak{a}_r \cap \mathfrak{b}_{t+1} \cap \cdots \cap
      \mathfrak{b}_s \not\subset p,
\]
we can choose an element $y \in P$ contained in the left-hand side and not
contained in $p$. Now $p_j' \subset p$ for $1 \le j \le t$, so that $y \notin p_1$, and hence
\[
  \mathfrak{a} : yA = \mathfrak{a}_1, \quad \text{and} \quad
  \mathfrak{b} : yA = \mathfrak{b}_1 \cap \cdots \cap \mathfrak{b}_t.
\]
Now set
\[
  B = A[x_1, \ldots, x_h], \quad \text{where } x_i = a_i/y,
\]
\[
  I = (x_1, \ldots, x_h)B \quad \text{and} \quad Q = PB + I.
\]

\textit{Step 2.} We prove that
\[
  B/I \cong A/\mathfrak{a}_1, \quad Q \in \Spec B \quad \text{and} \quad \htop Q = h + d.
\]
A general element of $B$ can be written in the form
\[
  a/y^\nu, \quad \text{with } a \in (\mathfrak{a} + yA)^\nu, \text{ for some } \nu \ge 0.
\]
Now $B = A + I$, so that $B/I \cong A/(I \cap A)$. Now if $\alpha \in I \cap A$, then there
exists $\nu > 0$ such that $y^\nu\alpha \in \mathfrak{a}(\mathfrak{a} + yA)^{\nu-1}$. Since
$\mathfrak{a} : yA = \mathfrak{a}_1$, we have $\alpha \in \mathfrak{a}_1$. Conversely,
$y\mathfrak{a}_1 \subset \mathfrak{a}$, giving $\mathfrak{a}_1 \subset I \cap A$, and hence $I \cap A = \mathfrak{a}_1$.
Therefore
\[
  B/I \cong A/\mathfrak{a}_1.
\]
Under this isomorphism, the prime ideals $p/\mathfrak{a}_1$ and $P/\mathfrak{a}_1$ of $A/\mathfrak{a}_1$
correspond to $(pB + I)/I$ and $(PB + I)/I$ respectively in $B/I$; hence, setting
\[
  \mathfrak{q} = pB + I = p + I \quad \text{and} \quad Q = PB + I = P + I,
\]
we have $\mathfrak{q}$, $Q \in \Spec B$, with $B/\mathfrak{q} = A/p$ and $B/Q = A/P$, and $Q$ is a
maximal ideal of $B$. Also, $B/I = A/\mathfrak{a}_1$, and since $I$ is generated by $h$
elements and $\mathfrak{a}_1$ is a $p$-primary ideal, we get
\[
  \htop Q = \dim B_Q \ge h + \htop(P/\mathfrak{a}_1) = h + \htop(P/p) = h + d.
\]
On the other hand, from $y \notin p$ we get $\mathfrak{q} = pA[y^{-1}] \cap B$: indeed, if
$a \in pA[y^{-1}] \cap B$, then we can write $a = c/y^\nu$ with
$c \in p \cap (\mathfrak{a} + y^\nu A)$; since $\mathfrak{a} \subset p$ and $y \notin p$, we have
\[
  p \cap (\mathfrak{a} + y^\nu A) = \mathfrak{a}(\mathfrak{a} + yA)^{\nu-1} + y^\nu p,
\]
and therefore $c \in I + p = \mathfrak{q}$. Also, $A[y^{-1}] = B[y^{-1}]$, so that
\[
  \htop \mathfrak{q} = \htop \mathfrak{q}A[y^{-1}] = \htop pA[y^{-1}] = \htop p = h.
\]
However, $Q/\mathfrak{q} = P/p$, so that
\[
  \htop(Q/\mathfrak{q}) = \htop(P/p) = d
\]
and therefore $\htop Q \ge d + h$. Putting this together with the previous
inequality, we see that $\htop Q = d + h$.

\textit{Step 3.} For $\nu = 1, 2, \ldots$, set
\[
  J_\nu = (x_1, \ldots, x_{h-1}, x_h - y^\nu)B,
\]
let $Q_\nu$ be a minimal prime divisor of $J_\nu$ satisfying
\[
  Q_\nu \subset Q \quad \text{and} \quad \htop(Q/Q_\nu) = \htop(Q/J_\nu),
\]
and set $P_\nu = Q_\nu \cap A$. We will complete the proof by showing that
$P_1, P_2, \ldots$ are all distinct, and that
\[
  \htop P_\nu = h, \quad \htop(P/P_\nu) = d \quad \text{for all } \nu.
\]
Now $B = A + J_\nu$, so that $B/Q_\nu \cong A/P_\nu$ and $Q/Q_\nu \cong P/P_\nu$, hence
$\htop(P/P_\nu) = \htop(Q/Q_\nu) = \htop(Q/J_\nu)$; moreover, $\htop Q = d + h$, and since
$J_\nu$ is generated by $h$ elements,
\[
  \htop(P/P_\nu) \ge d + h - h = d.
\]
Therefore we have
\[
  d \le \htop(P/P_\nu) = \htop(Q/Q_\nu) \le \htop Q - \htop Q_\nu = d + h - \htop Q_\nu,
\]
and so if we can prove that
\[
  (*) \quad \htop Q_\nu = \htop P_\nu \ge h
\]
then this will show simultaneously that $\htop(P/P_\nu) = d$ and $\htop P_\nu = h$.

We have already seen that $\mathfrak{q} = pA[y^{-1}] \cap B$, so that
$\mathfrak{q} \cap A = p$, and hence $y \notin \mathfrak{q}$. Also, $\mathfrak{a}_1$ is a $p$-primary ideal
with $B/I \cong A/\mathfrak{a}_1$, so that $I$ is a $\mathfrak{q}$-primary ideal, and from
$\htop \mathfrak{q} = h$ we get
\[
  \htop(I + yB) = \htop(x_1, \ldots, x_h, y)B = h + 1.
\]
In addition, we have $I + yB = J_\nu + yB$, and since $Q_\nu$ is a minimal prime
divisor of the ideal $J_\nu$, which is generated by $h$ elements,
$\htop Q_\nu \le h$, and hence $y \notin Q_\nu$, and
\[
  Q_\nu = Q_\nu A[y^{-1}] \cap B \quad \text{and} \quad P_\nu = Q_\nu A[y^{-1}] \cap A,
\]
so that, by \S\ref{sec:20}, Example 1,
\[
  \htop Q_\nu = \htop P_\nu.
\]
Furthermore
\[
  (\mathfrak{b} : yA)B \subset (x_1, \ldots, x_{h-1})B,
\]
\[
  (\mathfrak{b} : yA) + (a_h - y^{\nu+1})A \subset J_\nu \cap A \subset Q_\nu \cap A = P_\nu,
\]
and since all prime divisors of $(\mathfrak{b} : yA)$ are also prime divisors of
$\mathfrak{b}$, we have $\htop(\mathfrak{b} : yA) = h - 1$. Now $a_h \in \mathfrak{q}$ and
$y \notin p$, so that $a_h - y^{\nu+1} \notin p$, and since all minimal prime divisors of
$(\mathfrak{b} : yA)$ are contained in $p$ they do not contain $a_h - y^{\nu+1}$.
Therefore $\htop P_\nu \ge h$. This completes the proof of (*).

\textit{Step 4.} If $\nu < \mu$ then $(Q + Q_\mu)B_a$ contains $y^\nu = (y^\nu - y^\mu)/(1 - y^{\mu-\nu})$,
and therefore contains $(J_\nu + y^\nu B)B_a = (I + y^\nu B)B_a$, so that
\[
  \htop(Q + Q_\mu)B_a \ge h + 1, \quad \text{but } \htop Q_\nu \le h,
\]
and therefore $Q_\nu \ne Q_\mu$. From $Q_\nu A[y^{-1}] = P_\nu A[y^{-1}]$ we see that
$P_1, P_2, \ldots$ must also all be distinct.
\end{proof}

\begin{theorem}\label{thm:31.4}
A Noetherian local integral domain $(A, m)$ is catenary if and only if
\[
  \htop p + \coht p = \dim A \quad \text{for all } p \in \Spec A.
\]
\end{theorem}

\begin{proof}
`Only if' is trivial, and we prove `if'. Let $\dim A = n$; if $A$ is not catenary, then
there exist $p$, $P \in \Spec A$ such that
\[
  p \subset P, \quad \htop(P/p) = 1 \quad \text{but } \htop P > \htop p + 1.
\]
Set $\htop(m/P) = d$. Applying the strong existence theorem to $A/p$ we see that there
exist infinitely many $P_\lambda \in \Spec A$ such that
\[
  p \subset P_\lambda, \quad \htop(P_\lambda/p) = 1 \quad \text{and} \quad \htop(m/P_\lambda) = d.
\]
However, by assumption, $\htop(m/P_\lambda) + \htop P_\lambda = n$, so that
\[
  \htop P_\lambda = n - d = \htop P > \htop p + 1.
\]
But according to Theorem~\hyperref[thm:31.1]{1}, there are only finitely many such
$P_\lambda$, and we have a contradiction.
\end{proof}

If $A$ is a ring of finite Krull dimension, we say that $A$ is \emph{equidimensional} if
$\dim A/p = \dim A$ for every minimal prime $p$ of $A$.

\begin{lemma}\label{lem:31.2}
If an equidimensional local ring $(A, m)$ is catenary then
\[
  \htop p_2 = \htop p_1 + \htop(p_2/p_1) \quad \text{for all } p_1, p_2 \in \Spec A
      \text{ with } p_1 \subset p_2.
\]
\end{lemma}

\begin{proof}
If we choose a minimal prime ideal $p \subset p_1$ then
$\htop(p_2/p) = \htop(m/p) - \htop(m/p_2) = \dim A - \htop(m/p_2)$, and this is independent
of the choice of $p$, so that $\htop p_2 = \htop(p_2/p)$. Similarly,
$\htop p_1 = \htop(p_1/p)$, so that $\htop p_2 = \htop(p_2/p_1) + \htop p_1$.
\end{proof}

\begin{theorem}\label{thm:31.5}
Let $A$, $B$ be Noetherian local rings, and $A \to B$ a local homomorphism. If $B$ is
equidimensional and catenary and is flat over $A$ then $A$ is also
equidimensional and catenary, and $B/pB$ is equidimensional for every
$p \in \Spec A$.
\end{theorem}

\begin{proof}
Write $m$ and $\mathfrak{M}$ for the maximal ideals of $A$ and $B$. For any minimal prime ideal
$p_0$ of $A$ there exists a minimal prime ideal $P_0$ of $B$ lying over $p_0$; then
$\dim B/P_0 = \dim B$, so that $\dim B/p_0B = \dim B$, and then by
Theorem~\hyperref[thm:15.1]{15.1} we have
\[
  \htop(m/p_0) = \htop(\mathfrak{M}/p_0B) - \htop(\mathfrak{M}/mB) = \dim B - \htop(\mathfrak{M}/mB).
\]
This is independent of the choice of $p_0$, so that $A$ is equidimensional. If
$p \in \Spec A$ and $P$ is a minimal prime divisor of $pB$ then from the going-down
theorem (Theorem~\hyperref[thm:9.5]{9.5}) we see that $P \cap A = p$, so that by
Theorem~\hyperref[thm:15.1]{15.1}, $\htop P = \htop p$, and therefore
$\htop(\mathfrak{M}/P) = \htop \mathfrak{M} - \htop P = \htop \mathfrak{M} - \htop p$ is determined by $p$ only. That
is, $B/pB$ is equidimensional. Also, if $p' \in \Spec A$ is such that $p' \subset p$ and
$\htop(p/p') = 1$, we let $P'$ be a minimal prime divisor of $p'B$ contained in $P$;
then $B/p'B$ is also equidimensional and flat over $A/p'$, so that
$\htop(P/P') = \htop(P/p'B) = \htop(p/p') = 1$. However, $B$ is equidimensional and
catenary, so that $\htop(P/P') = \htop P - \htop P' = \htop p - \htop p'$, and therefore
$\htop p = \htop p' + 1$, and so $A$ is catenary.
\end{proof}

\begin{corollary*}
Let $A$ be a quotient of a regular local ring $R$. If $A$ is equidimensional then
so is its completion $A^*$.
\end{corollary*}

\begin{proof}
Let $P_0$ be a minimal prime ideal of $A^*$ and $p_0 = P_0 \cap A$; then writing
$\mathfrak{p} \subset R$ for the inverse image of $p_0$, we have
$R^*/\mathfrak{p}R^* = A^*/p_0A^*$. $R^*$ is an integral domain, and therefore
equidimensional, so we can apply the theorem to $R \to R^*$ and see that
$R^*/\mathfrak{p}R^*$ is equidimensional. Hence
\[
  \dim A^*/P_0 = \dim A^*/p_0A^* = \dim A/p_0 = \dim A.
\]
\end{proof}

\begin{definition*}
We say that a Noetherian local ring $A$ is \emph{formally equidimensional} (or
\emph{quasi-unmixed}) if its completion $A^*$ is equidimensional.
\end{definition*}

\begin{theorem}\label{thm:31.6}
Let $(A, m)$ be a formally equidimensional Noetherian local ring.
\begin{enumerate}[(i)]
  \item $A_p$ is formally equidimensional for every $p \in \Spec A$.
  \item If $I$ is an ideal of $A$, then
        \[
          A/I \text{ is equidimensional} \iff A/I \text{ is formally equidimensional}.
        \]
  \item If $B$ is a local ring which is essentially of finite type over $A$ (see
        p.~232) and if $B$ is equidimensional then it is also formally
        equidimensional.
  \item $A$ is universally catenary.
\end{enumerate}
\end{theorem}

\begin{proof}
(i) Let $P \in \Spec(A^*)$ be such that $P \cap A = p$, and set $B = (A^*)_P$.
Since $B$ is flat over $A_p$, by Theorem~\hyperref[thm:22.4]{22.4}, $B^*$ is flat over $(A_p)^*$. Now $B$ is
a quotient of a regular local ring, and is equidimensional, so that by the above
corollary, $B^*$ is also equidimensional. Hence by Theorem~\hyperref[thm:31.5]{5}, $(A_p)^*$ is also
equidimensional.

(ii) follows easily from Theorem~\hyperref[thm:31.5]{5}.

(iii) $B$ is a localisation of a quotient of $A[X_1, \ldots, X_n]$ for some $n$, so that
by (ii) we need only show that a localisation $B$ of $A[X_1, \ldots, X_n]$ is formally
equidimensional. Now $A^*[[X_1, \ldots, X_n]]$ is faithfully flat over
$A[X_1, \ldots, X_n]$, so that there is a local ring $C$ which is a localisation of
$A^*[[X_1, \ldots, X_n]]$, and a local homomorphism $B \to C$ such that $C$ is flat over
$B$, and hence $C^*$ is flat over $B^*$. By the remark after Theorem~\hyperref[thm:15.5]{15.5}, $C$ is
equidimensional, and is a quotient of a regular local ring, so that by the
corollary of Theorem~\hyperref[thm:31.5]{5}, $C^*$ is also equidimensional; hence by
Theorem~\hyperref[thm:31.5]{5}, $B^*$ is also equidimensional.

(iv) Any local integral domain essentially of finite type over $A$ is formally
equidimensional, and hence catenary, so that any integral domain which is
finitely generated over $A$ is catenary.
\end{proof}

We say that a Noetherian local ring $A$ is \emph{formally catenary} ([G2], (7.1.9)) if
$A/p$ is formally equidimensional for every $p \in \Spec A$. One sees easily from the
above theorem that formally catenary implies universally catenary. The converse
of this was proved by Ratliff [2]. Universally catenary is a property of finitely
generated $A$-algebras, and we have to deduce from this a property of the
completion, so that the proof is difficult. Before giving Ratliff's proof we make
the following observation.

Let $(R, m)$ be a Noetherian local integral domain, $K$ the field of fractions of $R$,
and $R'$ the integral closure of $R$ in $K$; let $S$ be an intermediate ring
$R \subset S \subset R'$ such that $S$ is a finite $R$-module. $S$ is a semilocal ring, and its
completion $S^*$ (with respect to the $m$-adic topology, which coincides with the
$\rad(S)$-adic topology) can be identified with $R^* \otimes_R S$. Now $R^*$ is flat over
$R$, so that $R \subset S \subset R' \subset K$ gives
$R^* \subset S^* \subset R^* \otimes_R R' \subset R^* \otimes_R K$. The ring $R^* \otimes_R K$ is the
localisation of $R^*$ with respect to $R - (0)$, so that writing $T$ for the total ring of
fractions of $R^*$, we can consider $R^* \otimes_R K \subset T$, and hence
$R^* \subset S^* \subset T$. This leads to the possibility that properties of $R^*$ will be
reflected in some $S$.

\begin{theorem}\label{thm:31.7}
For a Noetherian local ring $A$, the following conditions are equivalent:
\begin{enumerate}[(1)]
  \item $A$ is formally catenary,
  \item $A$ is universally catenary,
  \item $A[X]$ is catenary.
\end{enumerate}
\end{theorem}

\begin{proof}
(1) $\Rightarrow$ (2) was proved in Theorem~\hyperref[thm:31.6]{6} and (2) $\Rightarrow$ (3) is trivial. We
prove (3) $\Rightarrow$ (1). Suppose then that $A[X]$ is catenary; we will prove that $A^*$ is
equidimensional by assuming the contrary and deriving a contradiction. The
proof breaks up into several lemmas.

\begin{lemma}\label{lem:31.3}
Let $(R, m)$ be a catenary Noetherian local integral domain, and let $R^*$ be its
completion. Let $\dim R = n$, and suppose that there exists a minimal prime $Q$ of
$R^*$ such that
\[
  1 < \dim(R^*/Q) = d \le n.
\]
For $i = 1, 2, \ldots, d - 1$, write $\mathcal{O}_i$ for the set of $p \in \Spec R$ satisfying the
conditions
\begin{enumerate}[(1)]
  \item $\htop p = i$,
and (2) there exists a minimal prime divisor $P$ of $pR^*$ such that $Q \subset P$ and
$\dim(R^*/P) = d - i$.
\end{enumerate}
Then $\mathcal{O}_i$ is non-empty for each $i$.
\end{lemma}

\begin{proof}
We work by induction on $i$. If $0 \ne a \in m$ then any minimal prime divisor $P$ of
$aR^* + Q$ satisfies $\htop(P/Q) = 1$, $P \cap R \ne 0$ and contains $a$. Hence if we set
$\mathcal{M} = \{P \in \Spec(R^*) \mid Q \subset P, \htop(P/Q) = 1 \text{ and } P \cap R \ne 0\}$, then
$m = \bigcup_{P \in \mathcal{M}}(P \cap R)$. Now $\htop(mR^*/Q) = d > 1$, so that
$mR^* \notin \mathcal{M}$, and hence $m$ itself is not of the form $P \cap R$ for $P \in \mathcal{M}$;
therefore both $\mathcal{M}$ and $\{P \cap R \mid P \in \mathcal{M}\}$ are infinite sets. By
Theorem~\hyperref[thm:31.1]{1}, $\mathcal{M}' = \{P \in \mathcal{M} \mid \htop P = 1\}$ is also infinite; choose any
$P \in \mathcal{M}'$, and set $p = P \cap R$. Then $0 < \htop p = \htop pR^* \le \htop P = 1$, so that
$\htop p = 1$, and $P$ is a minimal prime divisor of $pR^*$. Since $R^*$ is catenary,
$\dim(R^*/P) = \dim(R^*/Q) - \htop(P/Q) = d - 1$; hence $p \in \mathcal{O}_1$ and the
assertion is true for $i = 1$.

If $i > 1$, take $p$ as above, and set $\bar{R} = R/p$ and $\bar{P} = P/pR^*$; then since
$\bar{R}$ is catenary, $\dim \bar{R} = n - 1$, and $\bar{P}$ is a minimal prime divisor of
$\bar{R}^* = R^*/pR^*$, with $\dim(\bar{R}^*/\bar{P}) = \dim(R^*/P) = d - 1 < n - 1$. Hence by
induction there exists a prime ideal $\bar{p}_i = p_i/p$ of $\bar{R}$ of height $i - 1$, and
a minimal prime divisor $\bar{P}_i$ of $\bar{p}_i\bar{R}^*$ such that $\bar{P} \subset \bar{P}_i$ and
$\dim(\bar{R}^*/\bar{P}_i) = (d - 1) - (i - 1) = d - i$. If $P_i = \bar{P}_i/pR^*$ then $P_i$ is a minimal
prime divisor of $p_iR^*$ containing $P$, and hence $Q$, and $R^*/P_i = \bar{R}^*/\bar{P}_i$ is
$(d - i)$-dimensional, so that from the fact that $R$ is catenary we get
$\htop p_i = \htop \bar{p}_i + \htop p = i$, and $p_i \in \mathcal{O}_i$.
\end{proof}

\begin{lemma}\label{lem:31.4}
Let $(R, m)$ be a Noetherian local integral domain, $R^*$ its completion, and let
$0 = \mathfrak{q}_1 \cap \cdots \cap \mathfrak{q}_r$ be a shortest primary decomposition of $0$ in $R^*$, with
$P_i = \sqrt{\mathfrak{q}_i}$ for $1 \le i \le r$. Suppose that $P_1$ satisfies
\[
  \htop P_1 = 0, \quad \coht P_1 = 1 < \dim R.
\]
Then there exist $b \in m$ and $\delta \in (\mathfrak{q}_2 \cap \cdots \cap \mathfrak{q}_r) - P_1$ with the
following properties:
\begin{enumerate}[(1)]
  \item $b - \delta \in \mathfrak{q}_1$,
  \item $(b, \delta)R^* = (b, c)R^*$,
and (3) $c/b \notin R$ but is integral over $R$.
\end{enumerate}
\end{lemma}

\begin{proof}
\textit{Step 1.} $P_1$ is a minimal prime ideal and $r > 1$, so that
$\mathfrak{q}_2 \cdots \mathfrak{q}_r \not\subset P_1$, and we can choose
$\delta' \in (\mathfrak{q}_2 \cap \cdots \cap \mathfrak{q}_r) - P_1$. Since $\coht P_1 = 1$, it follows that
$\mathfrak{q}_1 + \delta'R^*$ is $(mR^*)$-primary, and hence if we set
\[
  \mathfrak{a} = (\mathfrak{q}_1 + \delta'R^*) \cap R,
\]
then this is an $m$-primary ideal. Now if $0 \ne b \in \mathfrak{a}$ is any element, then in
$R^*$ we can write
\[
  b = \xi + \beta\delta' \quad \text{with } \xi \in \mathfrak{q}_1 \text{ and } \beta \in R^*.
\]
Since $b$ is not a zero-divisor in $R^*$ we have $\beta\delta' \notin P_1$, so that setting
$\delta = \beta\delta'$, we get (1).

\textit{Step 2.} We prove that $\delta \notin bR^*$. By contradiction, suppose that
$\delta = b\zeta$ with $\zeta \in R^*$; then $b - \delta = b(1 - \zeta) \in \mathfrak{q}_1$, and since
$b \notin P_1$ we have $1 - \zeta \in \mathfrak{q}_1 \subset mR^*$, so that $\zeta$ is a unit of $R^*$, and
$b\delta R^* \subset \mathfrak{q}_2$. This contradicts the fact that $b$ is not a zero-divisor of
$R^*$.

\textit{Step 3.} If $\mathfrak{b}$ is a non-zero ideal of $R$ such that $m$ is not a prime divisor of
$\mathfrak{b}$ then $\delta \notin \mathfrak{b}R^*$. Indeed, $\mathfrak{b} : m = \mathfrak{b}$, so that
$\mathfrak{b}R^* : mR^* = (\mathfrak{b} : m)R^* = \mathfrak{b}R^*$, and so $mR^*$ is not a prime divisor of
$\mathfrak{b}R^*$; if $P$ is any prime divisor of $\mathfrak{b}R^*$, then $P \ne mR^*$ and
$P \ne P_1$ (in view of $P_1 \cap R = 0$), so that $\coht P_1 = 1$ implies $P \not\supset P_1$,
hence $P \not\supset \mathfrak{q}_1$, and there exists $a \in \mathfrak{q}_1 - P$. If we write $\mathfrak{Q}$ for the
$P$-primary component of $\mathfrak{b}R^*$ then $a\delta = 0 \in \mathfrak{Q}$, so that
$\delta \in \mathfrak{Q}$. Thus finally $\delta \notin \mathfrak{b}R^*$.

\textit{Step 4.} By the previous two steps $m$ is a prime divisor of $bR$. Hence we can
write
\[
  bR = I \cap J \quad \text{with } I \text{ an } m\text{-primary ideal and } J : m = J,
\]
and then $bR^* = IR^* \cap JR^*$ with $\delta \in JR^*$ and $\delta \notin IR^*$. Moreover,
$(IR^* + JR^*)/IR^* \cong (I + J)/I$, so we can choose $c \in J$ such that
$\delta - c \in IR^*$. Then
\[
  \delta - c \in IR^* \cap JR^* = bR^*,
\]
so that $(b, c)R^* = (b, \delta)R^*$, and (2) is proved. If $c \in bR$ we would have
$(b, \delta)R^* = bR^*$, contradicting Step 2, so that $c/b \notin R$. On the other hand, we
have $b - \delta \in \mathfrak{q}_1$, so that $\delta(b - \delta) = 0$, that is $b\delta = \delta^2$, and
$c - \delta \in bR^*$. Set $c = \delta + b\eta$; then
$c^2 = \delta^2 + 2b\delta\eta + b^2\eta^2 \in b(\delta, b)R^* = b(c, b)R^*$, so that
\[
  c^2 \in (bc, b^2)R^* \cap R = (bc, b^2)R.
\]
From this, we get $c^2 = bcu + b^2v$ with $u, v \in R$, which proves that $c/b$ is integral
over $R$. Thus we have proved (3).
\end{proof}

\begin{lemma}\label{lem:31.5}
In the notation and assumptions of Lemma~\hyperref[lem:31.4]{4}, set $S = R[c/b]$; then $S$ has a
maximal ideal of height $1$.
\end{lemma}

\begin{proof}
Write $T$ for the total ring of fractions of $R^*$; then we can view $S^*$ as an
intermediate ring $R^* \subset S^* = R^*[c/b] \subset T$, and $T$ is the total ring of
fractions of $S^*$. In Lemma~\hyperref[lem:31.4]{4} we had $\Ass(R^*) = \{P_1, \ldots, P_r\}$, so that
setting $Q_i = P_iT \cap S^*$, we get
\[
  \Ass(S^*) = \{Q_1, \ldots, Q_r\} \quad \text{with } \htop Q_i = \htop P_i.
\]
Moreover, $S^*$ is integral over $R^*$, so that $S^*/Q_i$ is integral over $R^*/P_i$, and
hence also $\coht Q_i = \coht P_i$. Let $P^*$ be any maximal ideal of $S^*$ containing
$Q_1$. Then from $(b, c)R^* = (b, \delta)R^*$ we get $S^* = R^*[c/b] = R^*[\delta/b]$, and since
$\delta \in \mathfrak{q}_2 \cap \cdots \cap \mathfrak{q}_r$ and $\delta - b \in \mathfrak{q}_1$ we have
\[
  \delta/b \in Q_2 \cap \cdots \cap Q_r \quad \text{and} \quad \delta/b - 1 \in Q_1,
\]
so that $Q_i + Q_1 = S^*$ for all $i > 1$. Therefore $Q_1$ is the only minimal prime ideal
contained in $P^*$. However, $\coht Q_1 = 1$ and $P^* \cap R^* = mR^*$, so that
$\htop P^* = 1$. Setting $P = P^* \cap S$, we have $\htop P = \htop P^* = 1$, and $P$ is a maximal
ideal of $S$, since $P^*$ is a maximal ideal of $S^*$.
\end{proof}

Now we return to the proof of Theorem~\hyperref[thm:31.7]{7}. Let $A$ be a Noetherian local
integral domain, and suppose that $A^*$ is not equidimensional; then by
Lemma~\hyperref[lem:31.3]{3}, there exists a prime ideal $p$ of $A$ such that $(A/p)^* = A^*/pA^*$
has dimension $> 1$, and has a minimal prime ideal of coheight $1$. Set $R = A/p$;
then by Lemmas~\hyperref[lem:31.4]{4} and \hyperref[lem:31.5]{5} applied to $R$, there exists a subring $S$ of the
integral closure $R'$ of $R$ generated by one element, and having a maximal ideal
$P$ with $\htop P = 1 < \dim R$. Let $f\colon R[X] \to S$ be a surjective homomorphism of
$R$-algebras and let $\mathfrak{q}$ be its kernel. Set $Q = f^{-1}(P)$. Then $P = Q/\mathfrak{q}$ and
$Q \cap R = m$. Since $R \subset S$ we have $\mathfrak{q} \cap R = (0)$, hence
$\htop Q = \htop m + 1$ and $\htop \mathfrak{q} = 1$ by the remark after
Theorem~\hyperref[thm:15.5]{15.5}. Thus $\htop Q - \htop \mathfrak{q} = \htop m = \dim A/p > 1 = \htop(Q/\mathfrak{q})$,
so that $R[X]$ (and hence also $A[X]$) is not catenary.
\end{proof}

\begin{corollary*}
A Noetherian ring $A$ is universally catenary if and only if $A[X]$ is catenary.
\end{corollary*}

\begin{proof}
Suppose $A[X]$ is catenary and set $B = A[X_1, \ldots, X_n]$. In order to prove that
$B$ is catenary it suffices to prove that $B_P$ is catenary for every $P \in \Spec B$.
Let $p = P \cap A$. Then $B_P$ is a localisation of $A_p[X_1, \ldots, X_n]$. Since
$A_p[X_1]$ is catenary, $A_p[X_1, \ldots, X_n]$ is catenary by the theorem.
\end{proof}

\begin{corollary*}
A Noetherian ring of dimension $d$ is catenary if $d \le 2$ and is universally
catenary if $d \le 1$.
\end{corollary*}

\begin{proof}
The first assertion is obvious from the definitions and the second assertion
follows from the first because $\dim A[X] = d + 1$.
\end{proof}
